\setcounter{ex}{0}
\setcounter{paragraph}{0}
\PhanI
%\loai{Các biểu thức $\omega$, T, f}

\begin{ex} %[3P1Y3-1][Loai1]  
Tần số dao động của con lắc lò xo được tính theo biểu thức
\choice
{$f=\sqrt{\dfrac{m}{k}} $}
{$f=\sqrt{\dfrac{k}{m}} $}
{$f=\dfrac{1}{2\pi}\sqrt{\dfrac{m}{k}} $}
{\True $f=\dfrac{1}{2\pi}\sqrt{\dfrac{k}{m}} $}
\loigiai{
Tần số của dao động $f=\dfrac{1}{2\pi}\sqrt{\dfrac{k}{m}}$.
}
\end{ex} \haidong{20}
\begin{ex} 
(QG - 2022) Một con lắc lò xo gồm vật nhỏ khối lượng m và lò xo nhẹ có độ cứng k đang dao động điều hòa. Đại lượng $T=2\pi \sqrt{\dfrac{m}{k}}$ được gọi là
\choice
{biên độ dao động của con lắc}
{tần số của con lắc}
{tần số góc của con lắc}
{\True chu kì của con lắc}
\loigiai{
Chu kì của con lắc lò xo $T=2\pi \sqrt{\dfrac{m}{k}}$
}
\end{ex} \haidong{20}
%\loai{Sự thay đổi của T}

\begin{ex}%[3P1B3-1][Loai1]
Con lắc lò xo dao động điều hòa. Khi tăng khối lượng của vật lên 16 lần thì chu kì dao động của vật
\choice
{\True tăng lên 4 lần}
{giảm đi 4 lần}
{tăng lên 8 lần}
{giảm đi 8 lần}
\loigiai{
$T=2\pi \sqrt{\dfrac{m}{k}} \Rightarrow$ m tăng 16 lần thì T tăng lên 4 lần.
}
\end{ex} \haidong{20}
%\loai{Sự thay đổi tần số}


\begin{ex} %[3P1B3-1][Loai1]  
Khi treo vào con lắc lò xo có độ cứng $k_1$ một vật có khối lượng m thì vật dao động với chu kì $T_1$. Khi treo vật này vào lò xo có độ cứng $k_2$ thì vật dao động với chu kì $T_2 = 2T_1$ thì
\choice
{$k_1 = k_2$}
{\True $k_1 = 4k_2$}
{$k_2 = 2k_1$}
{$k_2 = 4k_1$}
\loigiai{
Ta có: $T=2\pi\sqrt{\dfrac{m}{k}} \Rightarrow k_1=4k_2$.
}
\end{ex} \haidong{20}

\begin{ex} %[3P1B3-1][Loai1]  
Trong con lắc lò xo nếu ta tăng khối lượng vật nặng lên 4 lần và độ cứng tăng 2 lần thì tần số dao động của vật
\choice
{tăng 2 lần}
{giảm 2 lần}
{tăng $\sqrt{2}$ lần}
{\True giảm $\sqrt{2}$ lần}
\loigiai{
\begin{itemize}
\item Ta có: $f_1 = \dfrac{1}{2\pi} \sqrt{\dfrac{k}{m}}$. Khi tăng m và k: $f_2 = \dfrac{1}{2\pi} \sqrt{\dfrac{2k}{4m}} = \dfrac{1}{2\pi} \sqrt{\dfrac{k}{m}}.\dfrac{1}{\sqrt{2}} = \dfrac{f_1}{\sqrt{2}}$.
\end{itemize}}

\end{ex} \haidong{20}


%\loai{Sự thay đổi của biên độ}

\begin{ex} %[3P1B3-1][Loai1] 
(QG - 2016) Một con lắc lò xo dao động điều hòa theo phương nằm ngang. Nếu biên độ dao động tăng gấp đôi thì tần số dao động điều hòa của con lắc
\choice
{tăng $\sqrt{2}$ lần}
{giảm 2 lần}
{\True không đổi}
{tăng 2 lần}
\loigiai{
Tần số dao động điều hòa của con lắc lò xo không phụ thuộc biên độ dao động.
}
\end{ex} \haidong{20}

%\loai{Loại khác}
\begin{ex} %[3P1B3-1][Loai1]  
Nếu m là khối lượng của vật, k là độ cứng của lò xo thì $\left(2\pi\sqrt{\dfrac{m}{k}}\right)$ có đơn vị là
\choice
{\True s (giây)}
{N (niutơn)}
{rad/s}
{Hz (héc)}
\loigiai{
Với m là khối lượng (kg), k là độ cứng của lò xo (N/m) thì $2\pi \sqrt{\dfrac{m}{k}}$ có đơn vị giây.
}
\end{ex} \haidong{20}

 
%\loai{Biểu thức}
\begin{ex} %[3P1Y5-1][Loai1] 
(QG - 2015) Một con lắc lò xo có khối lượng vật nhỏ là m dao động điều hòa theo phương ngang với phương trình $x = A\cos\omega t$. Mốc tính thế năng ở vị trí cân bằng. Cơ năng của con lắc là
\choice
{$m\omega A^2$}
{$\dfrac{1}{2}m\omega A^{2}$}
{$m\omega ^{2}A^{2}$}
{\True $\dfrac{1}{2}m\omega ^{2}A^{2}$}
\loigiai{
Cơ năng của vật dao động điều hòa: $W = W_{\text{đ}\max} = \dfrac{1}{2}mv_{\max}^{2} = \dfrac{1}{2}m\omega ^2A^2$.
}
\end{ex} \haidong{20}

\begin{ex}
(QG - 2021) Một con lắc lò xo gồm vật nhỏ và lò xo nhẹ có độ cứng k, đang dao động điều hòa theo phương nằm ngang. Mốc thế năng ở vị trí cân bằng. Khi vật qua vị trí có li độ x thì thế năng của con lắc là
\choice
{\True $W_t=\dfrac{1}{2}kx^2$}
{$W_t=kx$}
{$W_t=\dfrac{1}{2}kx$}
{$W_t=kx^2$}
\end{ex} \haidong{20}
\begin{ex} %[3P1Y5-1][Loai1] 
(QG - 2017) Một chất điểm có khối lượng m đang dao động điều hòa. Khi chất điểm có vận tốc v thì động năng của nó là
\choice
{\True $\dfrac{m{v}^{2}}{2}$}
{$mv^{2}$}
{$vm^{2}$}
{$\dfrac{v{m}^{2}}{2}$}
\loigiai{
Động năng của chất điểm được xác định bằng biểu thức $W_{\text{đ}}=\dfrac{1}{2}mv^{2}$.
}
\end{ex} \haidong{20}
%\loai{Cơ năng}
\begin{ex} %[3P1Y5-1][Loai1] 
(MH - 2017) Một con lắc lò xo đang dao động điều hòa, đại lượng nào sau đây của con lắc được bảo toàn?
\choice
{Cơ năng và thế năng}
{Động năng và thế năng}
{\True Cơ năng}
{Động năng}
\loigiai{
Cơ năng của con lắc được bảo toàn.
}
\end{ex} \haidong{20}
\begin{ex}
Cơ năng của một chất điểm dao động điều hoà tỉ lệ thuận với
\choice
{chu kì dao động}
{biên độ dao động}
{\True bình phương biên độ dao động}
{bình phương chu kì dao động}
\end{ex} \haidong{20}
\begin{ex} %[3P1Y5-1][Loai1]   
Năng lượng vật dao động điều hòa
\choice
{tỉ lệ với biên độ dao động}
{\True bằng thế năng của vật khi vật có li độ cực đại}
{bằng thế năng của vật khi vật qua vị trí cân bằng}
{bằng động năng của vật khi có li độ cực đại}
\loigiai{
Ta có: $W=\dfrac{1}{2}k.A^2$.
}
\end{ex} \haidong{20}

%\loai{Sự thay đổi của cơ năng}

\begin{ex}
Đại lượng nào sau đây tăng gấp đôi khi tăng gấp đôi biên độ của dao động điều hoà của con lắc lò xo?
\choice
{Cơ năng của con lắc}
{Động năng của con lắc}
{\True Vận tốc cực đại}
{Thế năng của con lắc}
\end{ex} \haidong{20}

%\loai{Chu kì, tần số của thế năng và động năng}
\begin{ex} %[3P1B5-1][Loai1]  
Một con lắc lò xo dao động điều hòa với tần số $2f_1$. Động năng của con lắc biến thiên tuần hoàn theo thời gian với tần số $f_2$ bằng
\choice
{$f_1$}
{$\dfrac{f_1}{2}$}
{\True $4f_1$}
{$2f_1$}
\loigiai{
\begin{itemize}
\item Phương trình dao động của vật là: $x=A.\cos \left({\omega t+\varphi }\right)$ $=A.\cos \left({2\pi ft+\varphi }\right)$.
\item Phương trình động năng:
\begin{align*}
W_{\text{đ}}=\dfrac{1}{4}.m.\omega ^2.A^2.\left({1-\cos \left({2\omega t+2\varphi }\right)}\right)=\dfrac{1}{4}.m.\omega ^2.A^2.\left({1-\cos \left({2.2\pi f.t+2\varphi }\right)}\right).
\end{align*}
\item Vậy tần số dao động của động năng bằng 2 lần tần số của dao động, nên khi tần số dao động tăng gấp đôi thì tần số của động năng tăng 4 lần $=4f_1$.
\end{itemize}
}%<MyLT1>
\end{ex} \haidong{20}


%\loai{Thế năng và động năng}

\begin{ex} %[3P1B5-1][Loai1] 
(QG - 2017) Một con lắc lò xo gồm vật nhỏ và lò xo nhẹ, đang dao động điều hòa trên mặt phẳng nằm ngang. Động năng của con lắc đạt giá trị cực tiểu khi
\choice
{lò xo không biến dạng}
{vật có vận tốc cực đại}
{vật đi qua vị trí cân bằng}
{\True lò xo có chiều dài cực đại}
\loigiai{
Động năng của vật đạt cực tiểu khi lò xo có chiều dài cực đại.
}
\end{ex} \haidong{20}

%\loai{Tổng hợp}

\begin{ex}%[3P1B5-1][Loai1]
Khi nói về dao động điều hoà của một chất điểm, phát biểu nào sau đây là \textbf{sai}?
\choice
{Khi động năng của chất điểm giảm thì thế năng của nó tăng}
{Biên độ dao động của chất điểm không đổi trong quá trình dao động}
{\True Độ lớn vận tốc của chất điểm tỉ lệ thuận với độ lớn li độ của nó}
{Cơ năng của chất điểm được bảo toàn}
\loigiai{
Phát biểu sai là “Độ lớn vận tốc của chất điểm tỉ lệ thuận với độ lớn li độ của nó”.
}
\end{ex} \haidong{20}
\begin{ex}
Phương trình dao động của một chất điểm dao động điều hoà là $x=A\cos(\omega t+\dfrac{2\pi}{3})$ cm. Động năng của nó biến thiên theo thời gian với biểu thức nào sau đây?
\choice
{$W_{\ct{đ}}=\dfrac{mA^2\omega^2}{2}\left[1+\cos(2\omega t+\dfrac{\pi}{3})\right]$}
{$W_{\ct{đ}}=\dfrac{mA^2\omega^2}{4}\left[1-\cos(2\omega t+\dfrac{4\pi}{3})\right]$}
{$W_{\ct{đ}}=\dfrac{mA^2\omega^2}{2}\left[1+\cos(2\omega t+\dfrac{4\pi}{3})\right]$}
{\True $W_{\ct{đ}}=\dfrac{mA^2\omega^2}{4}\left[1+\cos(2\omega t+\dfrac{4\pi}{3})\right]$}
\end{ex} \haidong{20}
\setcounter{ex}{0}
\PhanII
\begin{ex}
Cho các nhận định sau về động năng và thế năng của một vật dao động điều hòa. Chọn mốc thế năng tại vị trí cân bằng.
\choiceTF[t]
{Động năng đạt cực đại tại biên}
{\True Động năng và thế năng biến thiên với chu kì bằng một nửa chu kì dao động}
{\True Thế năng đạt cực tiểu tại vị trí cân bằng}
{Thế năng và động năng của vật biến thiên cùng tần số với tần số của li độ}
\loigiai{
\begin{itemchoice}
\itemch
\itemch
\itemch
\itemch
\end{itemchoice}

}
\end{ex} \haidong{20}

\begin{ex}
Cho các nhận định sau về động năng và thế năng của một vật dao động điều hòa. Chọn mốc thế năng tại vị trí cân bằng.
\choiceTF[t]
{\True Động năng đạt giá trị cực đại khi vectơ gia tốc đổi chiều}
{\True Thế năng đạt giá trị cực đại khi vật đổi chiều chuyển động}
{Động năng đạt giá trị cực tiểu khi vật đi qua vị trí cân bằng theo chiều âm của trục $O x$}
{\True Thế năng đạt cực tiểu khi vật chuyển từ chuyển động nhanh dần sang chậm dần}
\loigiai{
\begin{itemchoice}
\itemch
\itemch
\itemch
\itemch
\end{itemchoice}

}
\end{ex} \haidong{20}


\begin{ex}
Cho các nhận định sau về năng lượng trong dao động điều hòa của con lắc lò xo.
\choiceTF[t]
{\True Cơ năng của con lắc lò xo tỉ lệ với bình phương biên độ dao động}
{\True Có sự chuyển hóa qua lại giữa động năng và thế năng nhưng cơ năng được bảo toàn}
{\True Cơ năng của con lắc lò xo tỉ lệ với độ cứng $k$ của lò xo}
{Cơ năng của con lắc lò xo biến thiên theo quy luật hàm số sin với tần số bằng tần số của dao động điều hòa}
\loigiai{
\begin{itemchoice}
\itemch
\itemch
\itemch
\itemch
\end{itemchoice}

}
\end{ex} \haidong{20}

\begin{ex}
Cho các nhận định sau về đồ thị động năng thế năng của một vật dao động điều hòa.
\choiceTF[t]
{\True Đồ thị biểu diễn mối quan hệ giữa thế năng theo li độ là một phần của đường parabol}
{Đồ thị biểu diễn mối quan hệ giữa động năng theo thế năng là một đường hình sin}
{\True Đồ thị biểu diễn mối quan hệ giữa động năng theo li độ là một phần của đường parabol}
{Đồ thị biểu diễn mối quan hệ giữa động năng và tốc độ là một đoạn thẳng}
\loigiai{
\begin{itemchoice}
\itemch
\itemch
\itemch
\itemch
\end{itemchoice}

}
\end{ex} \haidong{20}




